We have a trained model, a judge, and a test set. The natural question
is the one the whole book has been building toward: \emph{does it
actually work?}

And the answer, as with almost everything in this book, is \textbf{``it
depends on what we mean,''} and then \textbf{``let's measure it rather
than argue.''} Because ``works'' is not a single number. It is a
question about \emph{which} kind of capability, under \emph{which}
conditions, measured in \emph{which} way. The job of this chapter is to
take that vague word and replace it with something we can actually see.

\section{The setup}\label{the-setup}

We trained a small decoder-only transformer on 1,400 generated tasks:
the model reads a natural-language specification and must produce the
reference Python program. The trained model is tiny --- \textbf{884,224
parameters} (0.88M), of which 53,632 are the embedding and 830,592 are
the actual computation. It saw 1,500 gradient steps on a batch of 16,
reaching a training loss of \textbf{0.0208}.

We then asked it to solve 300 \textbf{held-out} test tasks it never saw,
and scored each generated program by executing it against that task's
hidden tests. This is the moment of truth for the whole apparatus:
everything we built --- the world, the teacher, the judge, the rivals
--- converges on this measurement.

\section{What we measured: pass@1}\label{what-we-measured-pass1}

Here is the headline result, reported as \textbf{problem-level
hidden-test pass@1 under greedy decoding} --- that is, for each problem,
did the single program the model produced pass \emph{every} hidden test?

\begin{longtable}[]{@{}llll@{}}
\toprule\noalign{}
Candidate & syntax & executable & \textbf{pass@1} \\
\midrule\noalign{}
\endhead
\bottomrule\noalign{}
\endlastfoot
\textbf{MODEL} & 0.897 & 0.897 & \textbf{0.897} \\
TEMPLATE & 1.000 & 0.900 & 0.800 \\
RETRIEVAL & 1.000 & 1.000 & \textbf{1.000} \\
\end{longtable}

Three things stand out.

\textbf{First, the model genuinely learned something.} 0.897 means that
\textasciitilde90\% of the time, an unseen problem's natural-language
spec produced a program that passed all its hidden tests. That is not a
toy result. And crucially, it \textbf{beats the hand-written
template/rule baseline} (0.800). A hard-coded pattern-matching system
that recognizes keywords like ``sum'' and emits a template gets 80\%;
our learned model gets \textasciitilde90\%. So the model is doing
\emph{more} than a fixed rule system can do on this task.

\textbf{Second --- but retrieval does even better.} A trivial TF-IDF
retriever that does nothing but find the closest training example and
reuse its program scores \textbf{1.000}.

That second fact is not a failure; it is the single most important thing
this pilot taught us. It tells us the task, as currently constructed, is
\textbf{too easy to force generalization}. When a non-learning lookup
solves 100\% of the test set, then the test set is solvable by
retrieval, and a model scoring 90\% is not yet being tested on whether
it can generalize --- only on whether it can match the surface form of a
training example.

This is exactly the ``task too easy'' outcome the experimental plan
anticipated, and it is scientifically useful: it tells us \emph{how to
make the next task harder.}

\section{The moment that broke us (and the evaluator, not the
model)}\label{the-moment-that-broke-us-and-the-evaluator-not-the-model}

There is a story behind these numbers that is more instructive than the
numbers themselves, and it is worth telling plainly. When we ran the
evaluation for the very first time, the model scored \textbf{0.000}. Not
0.897 --- zero. Every single generated program failed the syntax check.

Our first instinct, naturally, was that the model was useless. We had
spent hours building a careful apparatus and the tiny model could not
even produce a parseable program. It would have been easy to conclude
``the model doesn't work'' and write the whole thing off.

But the book's own discipline --- \emph{test the tester} --- demanded we
ask whether the \emph{judge} was the problem, not the model. So we
looked at what the judge actually saw. And there it was: the generated
programs were perfectly valid Python, except they had a run of
\textbf{NUL bytes} (\texttt{\textbackslash{}x00}) appended after the
finished program. The model had learned to pad its output with a byte
that means ``nothing,'' and Python's parser rejects NUL bytes outright.
So \texttt{ast.parse} threw a syntax error on perfectly good code, and
the judge scored the model's valid program as a \emph{syntax failure}.

The fix was a single line: truncate the generated program at the first
NUL byte. A trivial change, and suddenly the model's score leapt from
0.000 to 0.897.

This is the single most important lesson of the whole book, and it is
not about the model at all. \textbf{The evaluator was broken, not the
model.} We had designed every experiment to detect shortcuts and
outright cheating in the \emph{model}, and we nearly missed the one case
where the measuring stick itself was faulty. The ``test the tester''
principle --- that we must verify the tool measures what we think it
measures before trusting a single number it produces --- was not a
formality. It just saved the entire project from a false negative.

The reason we caught it is worth pausing on. We did not catch it because
we are clever. We caught it because we had \emph{already} built a method
for testing testers (the mutation suite in Chapter 6), and we had
written the pilot's results to a manifest and looked at the intermediate
numbers, not just the headline. Had we grabbed the final pass@1 and
moved on, we would have shipped a book claiming a tiny model cannot
write Python --- a conclusion that was wrong.

\section{The shortcut audit}\label{the-shortcut-audit}

Before we trust the ``0.897'' number, we have to ask how the model might
have gotten there without learning.

\begin{enumerate}
\def\labelenumi{\arabic{enumi}.}
\item
  \textbf{Suspected shortcut: retrieval of a near-duplicate prompt.}
  \textbf{Why it could explain it:} the pilot's generated prompts are
  drawn from a small set of template families per operation, so many
  test prompts are surface-near to training prompts. \textbf{Control:}
  we ran the retrieval baseline --- 1.000. This \emph{confirms} the
  shortcut is available and that the test set is retrieval-solvable.
  \textbf{What it rules out:} nothing yet, because the model's 0.897 is
  \emph{below} the retrieval ceiling. It does not rule out the model
  partially relying on surface similarity. \textbf{What remains
  unresolved:} whether the model's own success is grounded in the
  \emph{computation} or in the \emph{surface form}. To answer that, we
  need the Level-3 (surface-change) and Level-4 (composition)
  experiments, which we run next.
\item
  \textbf{Suspected shortcut: exact memorization.} \textbf{Why it could
  explain it:} the model saw 1,400 examples; maybe it memorized them.
  \textbf{Control (Ch14):} on held-out prompts it still passed 59/60,
  and on a \emph{canary} prompt it produced the \emph{general} program
  rather than a memorized literal. \textbf{What it rules out:} pure
  example-memorization. \textbf{What remains:} the deeper question ---
  is the \emph{general program} it produces genuinely tied to the
  computation, or to the particular wording?
\end{enumerate}

\section{The honest summary}\label{the-honest-summary}

The pilot succeeded in exactly the way a pilot should: the laboratory
works, the model learns, and --- perhaps most valuably --- the
experiment told us something \emph{uncomfortable but true}. The task is
retrieval-solvable, and a model that beats a template rule system yet
sits below a lookup table has not yet been tested on the thing we care
about most: whether it generalizes.

And we were reminded, in the most direct way possible, that the
measuring stick is as much a subject of the experiment as the model. The
instrument and the thing measured both needed to be honest, and only one
of them initially was.

That is the baseline from which every later experiment is measured. From
here, we stop making the task easier and start making it \emph{probe}
generalization.
